\chapter{Introduction}\label{sec:introduction}

Relay communication uses an intermediate node to deliver information from a source to a destination. Classical amplify-and-forward (AF) and decode-and-forward (DF) provide well-studied baselines. Capacity results for coded block-DF apply under particular relay-channel assumptions and asymptotically long codes; they do not automatically characterize the uncoded symbol-wise DF studied here \cite{CoverGamal1979RelayChannel,ElGamalKim2011NetworkIT,Laneman2004CooperativeDiversity,RankovWittneben2007HalfDuplexRelay,Nosratinia2004CooperativeCommunication}.

Neural processing motivates a narrower experimental question: \emph{when does a small learned relay improve on the implemented classical alternatives, and what information and computation does that comparison require?} This thesis investigates that question by simulation, rather than proposing a universally superior relay or proving a capacity result.

Nine strategies are implemented, with eight evaluated in the canonical comparison: AF, symbol-wise DF, MLP, Hybrid, VAE, Transformer, Mamba-S6 and Mamba-2 SSD. A conditional GAN is implemented but excluded from that comparison on cost grounds. The source, channel models and destination are fixed within the canonical study; the learned models' architectures, training and checkpoints are part of the strategies being compared.

\paragraph{A sequence of bounded comparisons.} The canonical memoryless study establishes a reference under known receiver CSI. Separate fixed-ISI, composite-impairment, finite-pilot and blind studies then examine model mismatch and information limitations. These groups form a conceptual organization, not a one-factor-at-a-time ablation: modulation, fading, architecture and budgets differ where stated. A supplementary coded study investigates forwarding, soft read-out, rate selection and idealized readiness constraints.

\section{Scope and Contributions}\label{sec:system-model-intro}
\label{sec:introduction-scope}

The system has one half-duplex SISO relay and no direct source--destination path. Table~\ref{tbl:scope} distinguishes the principal experiments. Chapter~\ref{sec:methods} specifies their implementation and Chapter~\ref{sec:research-objectives} states the hypotheses.

\begin{longtable}{@{}p{0.24\linewidth}p{0.68\linewidth}@{}}
\caption[Scope of the thesis]{Primary experimental groups; conclusions apply to their stated configurations}\label{tbl:scope}\\
\toprule
Study & Configuration and measured question\\
\midrule
Canonical & Uncoded Gray-QPSK, i.i.d.\ per-symbol complex Rayleigh fading on both hops, receiver CSI; destination BER versus per-hop $E_s/N_0$ for eight relays.\\
Unknown/mismatched channel & Fixed BPSK ISI, faded-ISI QPSK, and specified composite families; model information, pilot budgets and blind comparators. Different groups are not causal controls for one another.\\
Coded supplement & Finite-length convolutional codes and a QPSK/16-QAM grid; implemented coded forwarding, rate selection, and idealized buffering. A coded-QPSK memory sweep also compares arithmetic and BER.\\
\bottomrule
\end{longtable}

In the canonical study, coherent equalization uses receiver CSI before the learned function is applied. No explicit channel coefficient is passed to that function. This is not a CSI-free chain. The unknown-channel MLPs are likewise trained on the evaluated impairment families, even when analytic parameters and test-time pilots are withheld.

The main contributions are:
\begin{enumerate}
\item A controlled canonical relay comparison, supplemented by equal-budget and replicated size studies. It distinguishes measured BER ordering from untested explanations involving capacity or overfitting.
\item Bounded experiments under ISI, nonlinearity and imperfect channel information, including matched Viterbi references, a pilot-count crossover at one SNR, and separately preserved fixed-budget rare-event validation.
\item A coded forwarding and memory/cost study that distinguishes symbol decisions, channel-sequence detection and error-correction decoding, while exposing the limitations of the implemented reliability metrics.
\end{enumerate}

The canonical DF measurements agree with the stated two-hop closed form within $3.6\%$ relative error at every tabulated SNR, after the QPSK $E_s/N_0$ conversion (Table~\ref{tbl:table2}). This is a simulator calibration check, not proof of all implementations or results. The additional controls and their endpoints are documented in their experiment sections.

\begin{longtable}{@{}p{0.24\linewidth}p{0.68\linewidth}@{}}
\caption[Assumptions and evidence boundaries]{Assumptions that limit interpretation of the results}\label{tbl:assumptions-claims}\\
\toprule
Dimension & Evidence boundary\\
\midrule
Processing & Symmetric learned windows use future samples; power normalization uses blocks. Neither this nor operation counts establishes causal streaming or measured end-to-end latency.\\
Channel knowledge & Canonical receiver CSI and training-family knowledge are available. Realizations vary within specified families; arbitrary family shift is untested.\\
Comparators & Viterbi MLSE optimizes sequence error, not necessarily bit error. QPSK BCJR is a relay-output control. Classical pipelines, including their metric mismatch, must be named explicitly.\\
Uncertainty & Fixed-checkpoint trials, replicated training, paired tests and nominal bit-level bounds have different units. Replication budgets are stated per study; null results do not prove equivalence and multiple comparisons are not jointly corrected.\\
Resources & Arithmetic and reference-program timing are implementation-specific. Parameter storage excludes runtime buffers and software. Energy, quantized hardware execution and deployment robustness are unmeasured.\\
\bottomrule
\end{longtable}

References to Viterbi in the coded supplement mean convolutional-code decoding; references to Viterbi MLSE in the ISI study mean channel-sequence detection. The combined coded-memory study uses both. The general relay channel with a direct link and cooperative-diversity results in the literature review are background, not the operating topology of this thesis.

\begin{figure}[htbp]
\noindent
\makebox[\textwidth][l]{%
\begin{tikzpicture}[
    >=latex,
    node distance=2.8cm,
    block/.style={
        rectangle, draw, thick,
        minimum width=2.3cm, minimum height=1.1cm,
        align=center
    },
    relayblock/.style={
        rectangle, draw, very thick, dashed,
        minimum width=2.4cm, minimum height=1.2cm,
        align=center
    }
]

\node[block] (S) {Source\\QPSK};
\node[relayblock, right=2.8cm of S] (R) {Relay\\$f_\theta(\cdot)$};
\node[block, right=2.8cm of R] (D) {Destination\\Decision};

\draw[->,thick] (S) -- node[above,align=center]
   {Hop 1\\$h_1[i],\,n_1[i]$} (R);

\draw[->,thick] (R) -- node[above,align=center]
   {Hop 2\\$h_2[i],\,n_2[i]$} (D);

% Note tied to the relay (the only varied block), no arrow into the datapath
\node[below=0.9cm of R, align=center, font=\footnotesize] (note)
   {$f_\theta(\cdot)\in\{$ AF, DF, MLP, Hybrid, VAE,\\
    Transformer, Mamba-S6, Mamba-2 SSD $\}$\\[2pt]
    \emph{(the only varied block; cGAN implemented but excluded here on cost grounds)}};

% thin dotted tie-line, clearly cosmetic, not a signal arrow
\draw[dotted] (R.south) -- (note.north);

\end{tikzpicture}%
}
\caption[Canonical two-hop relay benchmark]{Canonical two-hop relay benchmark used throughout the thesis. The source, the channel model on both hops, and the destination are the same in every experiment of the canonical comparison; only the relay processing function $f_\theta(\cdot)$ (dashed) is varied among the eight candidate strategies compared there. The fading itself is not static: $h_1[i]$ and $h_2[i]$ are redrawn every symbol. The comparison evaluates these implemented relay strategies, including their trained checkpoints}
\label{fig:canonical-relay-benchmark}
\end{figure}
